\documentclass[11pt,a4paper]{article}

\usepackage[utf8]{inputenc}
\usepackage[T1]{fontenc}
\usepackage[spanish]{babel}
\usepackage[margin=2cm, bottom=6.2cm, footskip=1.7cm]{geometry}
\usepackage{graphicx}
\usepackage{array}
\usepackage{needspace}
\graphicspath{{../../01_Propuesta/Logos/}{./}}
\usepackage{fancyhdr}
\usepackage{xcolor}
\usepackage[hidelinks]{hyperref}
\usepackage{enumitem}
\setlist{nosep}
\setlength{\parindent}{0pt}
\setlength{\parskip}{8pt}

\definecolor{verdeosc}{RGB}{26,107,60}
\definecolor{verdemid}{RGB}{46,139,87}

\pagestyle{fancy}
\fancyhf{}
\fancyfoot[C]{%
  \noindent{\color{verdeosc}\rule{\textwidth}{0.8pt}}\\[6pt]
  \tiny
  \hyphenpenalty=10000\exhyphenpenalty=10000
  \noindent
  \begin{minipage}{\textwidth}
  \begin{minipage}[t]{4.4cm}
  {\color{verdeosc}\textbf{ORGANIZA}}\\[3pt]
  Grupo de Investigación IN SILICO\\[6pt]
  {\color{verdeosc}\textbf{CO-ORGANIZADORES}}\\[3pt]
  Departamento de Biología y Química\\
  Departamento de Ingeniería Agroindustrial\\
  Departamento de Ingeniería Agrícola\\
  Departamento de Física\\
  Universidad de Sucre
  \end{minipage}%
  \hfill
  \begin{minipage}[t]{3.7cm}
  {\color{verdeosc}\textbf{COORDINACIÓN}}\\[3pt]
  \textbf{Aldo F. Combariza}\\
  Coordinación General\\
  \textbf{María Ximena Díaz}\\
  Coordinación Académica\\
  \textbf{Sebastián Vargas}\\
  Coordinación Logística\\
  \textbf{Selena Arias Avila}\\
  Coordinación Administrativa
  \end{minipage}%
  \hfill
  \begin{minipage}[t]{3.7cm}
  {\color{verdeosc}\textbf{CONTACTO}}\\[3pt]
  \textbf{SICVEC 2026}\\
  19--20 de octubre de 2026\\
  C.C. Guacarí, Sincelejo\\
  Modalidad híbrida\\
  \href{mailto:insilico@unisucre.edu.co}{insilico@unisucre.edu.co}
  \end{minipage}%
  \end{minipage}%
}
\renewcommand{\headrulewidth}{0pt}

\begin{document}

\noindent
\includegraphics[height=1.6cm]{logo_sicvec2026.png}%
\hspace{10pt}{\color{verdeosc}\vrule width 1.2pt height 1.5cm}\hspace{10pt}%
\includegraphics[height=1.28cm]{Logo_insilico.jpeg}%
\hspace{10pt}{\color{verdeosc}\vrule width 1.2pt height 1.5cm}\hspace{10pt}%
\includegraphics[height=1.3cm]{Logo_Bio.png}\hspace{6pt}%
\includegraphics[height=1.3cm]{logo_agroin.png}\hspace{6pt}%
\includegraphics[height=1.3cm]{logo_ingagrocola.png}\hspace{6pt}%
\includegraphics[height=1.3cm]{logo_fisica_lifi.png}%
\hspace{10pt}{\color{verdeosc}\vrule width 1.2pt height 1.5cm}\hspace{10pt}%
\includegraphics[height=1.3cm]{logo_unisucre.png}

\vspace{0.35cm}

\noindent
\textbf{Sincelejo, 8 de octubre de 2026}

\vspace{0.35cm}

\noindent
Señor\\
Guillermo Alejandro Corrales Ducuara\\
Estudiante, Maestría en Tecnologías Limpias\\
Universidad Santo Tomás, sede Bogotá D.C.\\
Bogotá D.C., Colombia

\vspace{0.3cm}

\noindent
\textbf{\large Asunto: Invitación como ponente — SICVEC 2026}

\vspace{0.3cm}

Estimado Guillermo:

\vspace{0.3cm}

En nombre del Grupo de Investigación IN SILICO y de los departamentos de Biología y Química, Ingeniería Agroindustrial, Ingeniería Agrícola y Física de la Universidad de Sucre (UNISUCRE), me complace invitarlo a participar como \textbf{ponente} en el \textbf{Simposio Internacional de Ciencia Verde y Economía Circular (SICVEC 2026)}.

Su trabajo de maestría sobre metano en codigestión anaeróbica, dentro de su formación en Tecnologías Limpias de la Universidad Santo Tomás, es justamente el tipo de investigación aplicada que buscamos en el eje de \textbf{Economía Circular}: la valorización de residuos orgánicos como fuente de energía. Nos gustaría que compartiera sus avances con investigadores, estudiantes y empresarios en el simposio.

\vspace{0.3cm}

\noindent
\textbf{INFORMACIÓN DEL EVENTO}

\renewcommand{\arraystretch}{1.3}
\begin{tabular}{ll}
\textbf{Nombre:} & Simposio Internacional de Ciencia Verde y Economía Circular \\
\textbf{Fechas:} & 19 y 20 de octubre de 2026 (lunes y martes) \\
\textbf{Lugar:} & Parque Comercial Guacarí, Sincelejo, Sucre, Colombia \\
\textbf{Modalidad:} & Híbrida \\
\textbf{Sitio web:} & \href{https://sicvec2026.eu.pythonanywhere.com}{sicvec2026.eu.pythonanywhere.com} \\
\end{tabular}

\vspace{0.3cm}

\needspace{8cm}
\noindent
\textbf{DETALLES DE LA PONENCIA}

\begin{tabular}{ll}
\textbf{Duración:} & 20 minutos (15 de charla + 5 de preguntas) \\
\textbf{Eje temático:} & Economía Circular \\
\textbf{Tema sugerido:} & Metano en codigestión anaeróbica \\
\textbf{Participación:} & Presencial, en el Parque Comercial Guacarí \\
\textbf{Franja:} & Martes 20 de octubre, por confirmar (entre las 10:00 y las 11:15 a.\,m.) \\
\textbf{Modalidad:} & Invitación académica \\
\end{tabular}

\vspace{0.3cm}

\needspace{4cm}
\noindent
\textbf{PRÓXIMOS PASOS}

\begin{enumerate}
  \item Confirmar su disponibilidad
  \item Proponer título definitivo y un breve resumen de su charla
  \item Enviar una reseña corta y una foto para la promoción del evento
\end{enumerate}

Rogamos responder esta invitación \textbf{lo antes posible} para incluir su charla en el programa.

\vspace{0.3cm}

Quedo atento a sus comentarios y a su disposición para coordinar cualquier detalle. Será un gusto contar con su participación en SICVEC 2026.

\vspace{0.35cm}

Atentamente,

\vspace{0.35cm}

\textbf{Dr. Aldo F. Combariza}\\
Coordinador General, SICVEC 2026\\
Grupo de Investigación IN SILICO · Departamento de Biología y Química\\
Facultad de Educación y Ciencias · Universidad de Sucre\\
\href{mailto:aldo.combariza@unisucre.edu.co}{aldo.combariza@unisucre.edu.co}

\vspace{0.2cm}

\end{document}
